\begin{helper}
Fix the history cell \(H\), terminal coordinate set \(U\), blue support
\(B\), red label \(J\), branch labels \(\beta\), and transcript labels
\(\tau=(P_\tau,A_\tau,Z_\tau)\) as in the red residual leaf construction.
Assume the terminal product law on \(U\), \(0\le p\le 1/2\), the fixed blue
support hypotheses for \(B\), the branch functionality from complete blue
terminal traces, the realized-cell geometry for \(R_{\beta,J}\), \(P_\tau\),
\(A_\tau\), and \(Z_\tau\), the open-candidate support exhaustion, the
selected-sibling pruning clause, the enlarged-cylinder span and disjointness
clauses, and the forward, reverse, and functional compatibility clauses.

Then there is a nonnegative mass
\[
\nu:\operatorname{Option}(\mathcal T\times\mathcal P)\to\mathbb R
\]
such that
\[
\sum_{\lambda\in \operatorname{Option}(\mathcal T\times\mathcal P)}
  \nu(\lambda)\le 1,
\]
and for every realized \((\beta,\tau)\),
\[
p^{|P_\tau|}(1-p)^{|A_\tau|}
\le
p^{u_R}p^{u_B}(1-p)^{|Z_\tau|}\nu(\operatorname{some}(\beta,\tau)).
\]
The factors for absent blue coordinates in the complete branch trace remain
inside \(\nu\); only the fixed present factors from \(B\) and
\(R_{\beta,J}\), together with the released selected absences \(Z_\tau\), are
factored out.
\end{helper}
\begin{proof}
Write \(q=1-p\).  For a realized pair \((\beta,\tau)\), define the enlarged
released cylinder \(C_{\beta,\tau}\) to be the family of complete terminal
assignments \(A\subseteq U\) satisfying the right-hand side of the span
hypothesis: \(P_\tau\subseteq A\), \(A\) is disjoint from
\(A_\tau\setminus Z_\tau\), and on every blue terminal coordinate outside
\(Z_\tau\) the membership of \(A\) agrees with
\(\operatorname{blueTrace}(\beta)\).  Since the realized-cell geometry says
that every member of \(Z_\tau\) is red, no absent blue-coordinate condition is
released.  The span clause identifies this cylinder with the disjoint union
of the compatible assignments for \((\beta,\tau)\) and the selected present
sibling pieces.  The selected-sibling pruning clause says that sibling pieces
are not compatible with any realized cell.  The disjointness clause says that
the cylinders \(C_{\beta,\tau}\) for distinct realized pairs are pairwise
disjoint; equivalently, by
\(\noderef{FixedSetHistoryCellRedFunctionalScanDisjointness}\), a complete
terminal assignment cannot have two realized compatible branch/transcript
outputs.

By the terminal product law supplied by
\(\noderef{FixedSetHistoryCellTerminalProductPartition}\), the conditional
mass of \(C_{\beta,\tau}\) is the product mass obtained by forcing the
coordinates of \(P_\tau\) present, forcing the coordinates of
\(A_\tau\setminus Z_\tau\) absent, and leaving all other terminal coordinates
free.  Thus
\[
\mu(C_{\beta,\tau})
=p^{|P_\tau|}q^{|A_\tau|-|Z_\tau|}.
\]
The containment \(B\cup R_{\beta,J}\subseteq P_\tau\), the cardinalities
\(|B|=u_B\) and \(|R_{\beta,J}|=u_R\), the blue-only and red-only color
conditions, and the disjointness of present and absent transcript coordinates
allow the fixed present factors \(p^{u_B}\) and \(p^{u_R}\) to be factored
from this cylinder mass.  The red support is determined once the blue trace
is fixed, by
\(\noderef{FixedSetHistoryCellRedSupportDeterminedByBlueTrace}\), so the
events for different branch traces do not introduce extra multiplicity from
several red-support choices.

The finite released-cylinder union estimate is isolated in
\(\noderef{FixedSetHistoryCellRedReleasedCylinderWeightNormalization}\).
Applied to the present data, it gives
\[
\sum_{\beta,\tau:\mathsf{Realized}(\beta,\tau)}
  p^{|P_\tau|}q^{|A_\tau|-|Z_\tau|}
\le p^{u_R}p^{u_B}.
\]
All absent blue-coordinate factors are still part of this released-cylinder
weight; no reciprocal absent-blue factor is introduced.

Define \(\nu(\operatorname{none})=0\), define
\(\nu(\operatorname{some}(\beta,\tau))=0\) for non-realized pairs, and for a
realized pair proceed as follows.  If \(p^{u_R}p^{u_B}>0\), define
\[
\nu(\operatorname{some}(\beta,\tau))
=
\frac{\mu(C_{\beta,\tau})}{p^{u_R}p^{u_B}}.
\]
If \(p^{u_R}p^{u_B}=0\), define
\(\nu(\operatorname{some}(\beta,\tau))=0\).  This is the only degenerate
case requiring a convention.

Nonnegativity follows from \(0\le p\le 1/2\), hence \(0\le q\le 1\), and
from nonnegativity of finite product masses.  In the nondegenerate case,
\(\noderef{FixedSetHistoryCellRedReleasedCylinderWeightNormalization}\)
gives
\[
\sum_{\beta,\tau:\mathsf{Realized}(\beta,\tau)}
  \mu(C_{\beta,\tau})
\le p^{u_R}p^{u_B}.
\]
Therefore
\[
\sum_{\beta,\tau:\mathsf{Realized}(\beta,\tau)}
  \nu(\operatorname{some}(\beta,\tau))
=
\frac{\sum_{\mathsf{Realized}(\beta,\tau)}\mu(C_{\beta,\tau})}
     {p^{u_R}p^{u_B}}
\le 1,
\]
using the positive fixed denominator.  In the degenerate case all assigned
masses are zero, so the same inequality is immediate.  The
\(\operatorname{none}\) and non-realized terms are zero in both cases, so
the displayed bound is the full sum over
\(\operatorname{Option}(\mathcal T\times\mathcal P)\).

It remains to check the domination inequality for a realized pair.  In the
degenerate case \(p^{u_R}p^{u_B}=0\), the finite union estimate from
\(\noderef{FixedSetHistoryCellRedReleasedCylinderWeightNormalization}\)
says that a sum of nonnegative released weights is at most \(0\).  Hence the
released weight
\(p^{|P_\tau|}q^{|A_\tau|-|Z_\tau|}\) of the realized pair is \(0\), and
multiplying by the nonnegative factor \(q^{|Z_\tau|}\) gives
\(p^{|P_\tau|}q^{|A_\tau|}=0\).  The right side is also \(0\) by the
definition of \(\nu\), so the inequality holds.  In the nondegenerate case,
multiplying the definition of \(\nu\) by \(p^{u_R}p^{u_B}\) gives
\[
p^{u_R}p^{u_B}\nu(\operatorname{some}(\beta,\tau))
=\mu(C_{\beta,\tau})
=p^{|P_\tau|}q^{|A_\tau|-|Z_\tau|}.
\]
Multiplying both sides by \(q^{|Z_\tau|}\) gives
\[
p^{|P_\tau|}q^{|A_\tau|}
=
p^{u_R}p^{u_B}q^{|Z_\tau|}
  \nu(\operatorname{some}(\beta,\tau)).
\]
The equality uses \(Z_\tau\subseteq A_\tau\), so
\(|A_\tau|=(|A_\tau|-|Z_\tau|)+|Z_\tau|\) in the exponent calculation.  This
is the required domination inequality, and the normalization has divided only
by the fixed present factors \(p^{u_R}p^{u_B}\).
\end{proof}
